\section{从经典几何到学习化建模的过渡}

\subsection{传统多视图几何重建范式}

传统多视图几何重建通常建立在 Structure-from-Motion（SfM）与 Multi-View Stereo（MVS）框架之上，其核心流程包括局部特征提取、跨图像匹配、几何验证、相机位姿恢复、三角化以及后续的全局优化\cite{triggs2000bundle,schoenberger2016colmap}。这一范式的优势在于几何解释清晰、优化目标明确，但同时也高度依赖局部描述子的稳定性与初始化质量。只要前端匹配失败，后续相机估计与重建质量往往会快速恶化。

经典管线的另一个特点是模块化程度高。前端负责提取与关联二维观测，后端则负责全局一致性优化。这种分治结构在工程实现上较为成熟，但也使得整个系统难以直接从最终重建误差中反向修正前端表征\cite{wang2024vggsfm}。随着深度学习方法在视觉任务中的广泛成功，研究者开始思考：能否让神经网络承担传统管线中最脆弱或最耗时的部分，同时保留多视图几何的结构性约束。

\subsection{学习方法为何进入多视图几何重建}

学习方法之所以逐步进入多视图几何重建，根本原因在于传统方法在复杂真实场景中的鲁棒性和泛化能力存在明显天花板。面对大视角变化、弱纹理、动态区域与复杂光照时，纯几何启发式往往难以稳定建立高质量对应\cite{wang2021deep2view}。而深度网络则能够从大规模数据中学习跨视角不变特征，并利用上下文提高匹配与估计质量。

不过，多视图几何重建的学习化并不意味着完全放弃几何结构。相反，该方向真正有效的工作通常强调“以几何为骨架、以学习为增强”的思路：学习模块负责估计更可靠的特征、轨迹或初始化，而多视图几何仍然提供问题定义、约束形式与评价标准\cite{bhalgat2023lighttouch,wang2024vggsfm}。从这一角度看，学习方法在该领域的发展更接近“重新组织几何管线”，而非“黑盒替代几何”。

\subsection{VGG 团队早期代表工作}

VGG 团队的早期工作体现出鲜明的过渡特征。首先，\textit{Unsupervised Learning of Probably Symmetric Deformable 3D Objects from Images in the Wild} 关注的是在弱监督条件下利用对称性等物理先验恢复三维结构，体现了“先验驱动学习”的思路\cite{wu2020unsup3d}。该工作虽然更接近单图重建，但它说明了几何与形状先验在学习式三维理解中的基础作用。

其次，\textit{Deep Two-View Structure-from-Motion Revisited} 则更加直接地讨论双视图结构恢复中哪些环节适合学习、哪些环节应保留经典几何求解\cite{wang2021deep2view}。该工作并未简单回归相机位姿或深度，而是将光流估计、相对位姿恢复和尺度不变深度估计有机组合，强调经典几何问题的可解性边界。这类工作为后续更系统的可微 SfM 研究奠定了方法论基础。

再次，\textit{Common Objects in 3D}（CO3D）数据集工作则从数据层面补齐了真实世界多视图学习的关键基础设施\cite{reizenstein2021co3d}。由于许多学习方法长期依赖合成数据或小规模真实数据，跨域泛化能力受限，CO3D 提供的大规模真实多视图对象数据集为后续研究提供了更接近真实场景的训练与评测基准。

\subsection{本综述的分析视角}

基于上述背景，本文后续不按论文年份简单罗列，而是按照学习方法逐步介入传统多视图几何管线的方式来组织。具体而言，本文将从“几何对应学习”“可微结构恢复”“统一前馈预测”三个层面梳理 VGG 团队的代表工作，并在必要处引入 DUSt3R 等相关方法作为横向对照\cite{wang2024dust3r}。
